
\loai{Khối lượng còn lại}
\begin{vidu}%[3P7-7.3Y][Loai4]
Ban đầu có 100 g phóng xạ với chu kì bán rã là 7 ngày đêm. Sau 28 ngày đêm khối lượng chất phóng xạ còn lại là
\choice
{93,75 g}
{87,5 g}
{12,5 g}
{\True 6,25 g}
\loigiai{
Khối lượng chất phóng xạ còn lại là:
\begin{align*}
m=m_0.2^{-\tfrac{t}{T}}=100.2^{-\tfrac{28}{7}}=6,25\ g.
\end{align*}}
\end{vidu}
\loai{Tính khối lượng ban đầu}
\begin{vidu}%[3P7-7.3B][Loai5]
Ban đầu một mẫu chất phóng xạ nguyên chất có khối lượng $m_0$, chu kì bán rã của chất này là 3,8 ngày. Sau 15,2 ngày khối lượng của chất phóng xạ đó còn lại là 2,24 g. Khối lượng $m_0$ là
\choice
{5,60 g}
{\True 35,84 g}
{17,92 g}
{8,96 g}
\loigiai{
Khối lượng còn lại sau khoảng thời gian t: 
\begin{align*}
m_t=m_0. 2^{ - \tfrac{t}{T}}\Rightarrow 2,24=m_0{.2^{ - \tfrac{15,2}{3,8}}}\Rightarrow m_0=35,84\ g.
\end{align*}
}
\end{vidu}
\loai{Số nguyên tử còn lại}
\begin{vidu}%[3P7-7.3B][Loai6]
Hạt nhân polonium $\Hn{84}{210}{Po}$ là chất phóng xạ có chu kì bán rã 138 ngày. Khối lượng ban đầu là 10 g. Cho $N_A = 6,023.10^{23}\ mol^{-1}$.  Lấy khối lượng nguyên tử tính theo đơn vị u bằng số khối của hạt nhân của nguyên tử đó. Số nguyên tử còn lại sau 207 ngày là
\choice
{$1,01.10^{23}$}
{\True $1,01.10^{22}$}
{$2,05.10^{22}$}
{$3,02.10^{22}$}
\loigiai{
\begin{itemize}
\item Ta có số hạt nhân Pôlôni ban đầu: $N_{0}=\dfrac{m_0}{M}.N_A=2,868.10^{22}$ nguyên tử.
\item Số nguyên tử còn lại sau 207 ngày là: 
\begin{align*}
N=N_{0}.2^{- \tfrac{207}{138}}=1,01.10^{22}\ \text{nguyên tử}.
\end{align*}
\end{itemize}
}
\end{vidu}
\loai{Tính thời gian}
\begin{vidu}%[3P7-7.3B][Loai7]
Chất phóng xạ $^{210}Po$ phát ra tia $\alpha$ và biến đổi thành $\Hn{82}{206}{Pb}$. Chu kì bán rã của $^{210}Po$  là 138 ngày. Ban đầu có 100 g Po thì sau bao lâu lượng Po chỉ còn 1 g?
\choice
{\True 916,85 ngày}
{834,45 ngày}
{653,28 ngày}
{548,69 ngày}	
\loigiai{
Khối lượng chất phóng xạ $^{210}Po$ còn lại:
\begin{align*}
m=m_0.2^{-\tfrac{t}{T}} \Rightarrow \dfrac{m}{m_0}=2^{ - \tfrac{t}{T}}=\dfrac{1}{100}\Leftrightarrow\dfrac{t}{138}=6,64\Leftrightarrow t=916,85\ \text{ngày}.
\end{align*}	
}
\end{vidu}
\loai{Phần trăm lượng còn lại và ban đầu}
\begin{vidu}%[3P7-7.3B][Loai8]
Gọi t là khoảng thời gian để số hạt nhân của một đồng vị phóng xạ giảm đi 4 lần. Sau thời gian 2t hạt nhân còn lại của đồng vị đó bằng bao nhiêu phần trăm số hạt nhân ban đầu 
\choice
{$25,25\%$}
{$93,75\%$}
{\True $6,25\%$}
{$13,5\%$}
\loigiai{
Ta có: $\left\{\begin{aligned}& \dfrac{N}{N_0}=\dfrac{1}{4}=2^{-\tfrac{t}{T}} \\
& \dfrac{N'}{N_0}=2^{-\tfrac{2t}{T}}=\left({2^{-\tfrac{t}{T}}}\right)^2 
\end{aligned}\right.\Rightarrow \dfrac{N'}{N_0}=\dfrac{1}{16}=6,25\%$.
}
\end{vidu}
\loai{Tính chu kì bán rã}
\begin{vidu}%[3P7-7.3B][Loai9]
Giả sử sau 3 giờ phóng xạ, số hạt nhân của một đồng vị phóng xạ còn lại bằng $25\%$ số hạt nhân ban đầu thì chu kì bán rã của đồng vị đó bằng
\choice
{2 giờ}
{1 giờ}
{\True 1,5 giờ}
{0,5 giờ}
\loigiai{
\begin{itemize}
\item Công thức tính số hạt nhân còn lại sau khoảng thời gian t: $N=N_0.2^{-\tfrac{t}{T}} $.
\item Sau 3 giờ: $\dfrac{N}{N_0}=2^{-\tfrac{t}{T}}\Leftrightarrow \dfrac{1}{4}=2^{-\tfrac{3}{T}}\Leftrightarrow T=1,5$ giờ.
\end{itemize}
}
\end{vidu}
\loai{Hai chất phóng xạ}
\begin{vidu}%[3P7-7.3B][Loai10]
Chất phóng xạ X có chu kì $T_1$, chất phóng xạ Y có chu kì $T_2=0,5T_1 $. Sau khoảng thời gian $t=T_1 $, thì khối lượng của chất phóng xạ còn lại so với khối lượng lúc đầu là
\choice
{\True X còn $\dfrac{1}{2}$; Y còn $\dfrac{1}{4}$}
{X còn $\dfrac{1}{4}$, Y còn $\dfrac{1}{2}$}
{X và Y đều còn $\dfrac{1}{4}$}
{X và Y đều còn $\dfrac{1}{2}$}
\loigiai{
\begin{itemize}
\item Số hạt nhân X còn lại sau $t=T_1$: $N_X=N_{0X}.2^{ - \tfrac{T_1}{T_1}}=\dfrac{N_{0X}}{2}$.
\item Số hạt nhân Y còn lại sau $t=T_1$: $N_Y=N_{0Y}.2^{-\tfrac{T_1}{0,5T_1}}=\dfrac{N_{0Y}}{4}$.
\end{itemize}
}
\end{vidu}
\begin{vidu}%[3P7-7.3K][Loai10]
Có hai chất phóng xạ A và B với hằng số phóng xạ $\lambda_A$ và $\lambda_B$. Số hạt nhân ban đầu trong 2 chất là $N_A$ và $N_B$. Thời gian để số hạt nhân A và B của hai chất còn lại bằng nhau là
\choice
{$\dfrac{\lambda_A\lambda_B}{\lambda_A-\lambda_B}\ln \dfrac{N_A}{N_B} $}
{$\dfrac{1}{\lambda_A+\lambda_B}\ln \dfrac{N_B}{N_A}$}
{\True $\dfrac{1}{\lambda_B-\lambda_A}\ln \dfrac{N_B}{N_A}$}
{$\dfrac{\lambda_A\lambda_B}{\lambda_A+\lambda_B}\ln \dfrac{N_A}{N_B}$}
\loigiai{
\begin{itemize}
\item Số hạt nhân A còn lại sau $t$: $N_{A1}=N_A.e^{-\lambda_A t}\quad (1)$.
\item Số hạt nhân B còn lại sau $t$: $N_{B1}=N_B.e^{-\lambda_B t}\quad (2)$.
\item Lập tỉ số $(1)$ và $(2)$: 
\begin{align*}
\dfrac{N_{A1}}{N_{B1}}=\dfrac{{N_A.e^{-\lambda _At}}}{{N_B.e^{-\lambda _Bt}}}=1
&\Leftrightarrow e^{-(\lambda _A-\lambda _B)t}=\dfrac{{N_B}}{{N_A}}
\Leftrightarrow -(\lambda _A-\lambda _B)t=ln\dfrac{{N_B}}{{N_A}}\\
&\Leftrightarrow t=-\dfrac{1}{{\lambda _A-\lambda _B}}\ln \dfrac{{N_B}}{{N_A}}=\dfrac{1}{{\lambda_B-\lambda_A}}\ln \dfrac{{N_B}}{{N_A}}.
\end{align*}
\end{itemize}
}
\end{vidu}
\loai{Bài toán tổng hợp}
\begin{vidu}%[3P7-7.3G][Loai14]
	\nguon{ĐH - 2013} Hiện nay uranium tự nhiên chứa hai đồng vị phóng xạ $^{235}\mathrm {U}$ và $^ {238} \mathrm {U}$, với tỉ lệ số hạt $^{235}{U}$ và số hạt $^{238}{U}$  là $\dfrac{7}{1000}$. Biết chu kì bán rã của  $^{235}{U}$ và $^{238}U$ lần lượt là $7,00.10^8$ năm và $4,50.10^9$ năm. Cách đây bao nhiêu năm, uranium tự nhiên có tỉ lệ số hạt $^{235}U$ và số hạt $^{238}U$ là $\dfrac{3}{100}$?
	\choice
	{2,74 tỉ năm}
	{2,22 tỉ năm}
	{\True 1,74 tỉ năm}
	{3,15 tỉ năm}
	\loigiai{
		\begin{itemize}
			\item Số hạt nhân $^{235}U$ còn lại hiện nay: $N_1=N_{01}.2^{-\tfrac{t}{T_1}}$.
			\item Số hạt nhân $^{238}U$ còn lại hiện nay: $N_2=N_{02}.2^{-\tfrac{t}{T_2}}$.
			\item Theo bài ra ta có: 
			\begin{align*}
				\begin{cases}
					\dfrac{N_{01}}{N_{02}}=\dfrac{3}{100}\\
					\dfrac{N_1}{N_2}=\dfrac{7}{1000}	
				\end{cases} \Rightarrow \dfrac{7}{1000}=\dfrac{3}{100}.2^{\tfrac{t}{T_2}-\tfrac{t}{T_1}} \Rightarrow t =\dfrac{\ln\dfrac{7}{30}}{\left(\dfrac{1}{T_2} - \dfrac{1}{T_2}\right)\ln 2}= 1,74.10^9\ \text{năm}.
			\end{align*}
		\end{itemize}
	}
\end{vidu}






















